\documentclass[../../../include/open-logic-section]{subfiles}

\begin{document}

\olfileid{sth}{ord-arithmetic}{add}
\olsection{ترتیبی جوڑ}

منّ لو اسیں $\alpha$ تے $\beta$ نوں جوڑنا چاہندے آں۔ اسیں $\beta$ دی اک {نقل}، $\alpha$ دی اک نقل توں عین پچھوں رکھ سکدے آں۔ (سانوں \emph{نقلاں} لینیاں پیندیاں نیں، کیوں جے \olref[ordinals][basic]{ordinalsaresubsets} توں سانوں پتا اے کہ یا تاں $\alpha
\subseteq \beta$ یا $\beta \subseteq \alpha$۔) سہج سمجھ مطابق ایہدا اثر ایہ اے کہ اسیں قدماں دی اک $\alpha$ لڑی توں لنگھدے آں \emph{تے فیر} اک $\beta$ لڑی توں لنگھدے آں۔ جیہڑی لڑی بنے گی اوہ حُسنِ ترتیب والی ہووے گی۔ سو \olref[ordinals][ordtype]{thmOrdinalRepresentation} مطابق اوہ اک (اکّو اک) ترتیبی عدد نال ہم ساخت اے۔ اوس ترتیبی عدد نوں $\alpha$ تے $\beta$ دا \emph{جوڑ} سمجھ سکدے آں۔

ایہ ترتیبی جوڑ دے پچھے سہج خیال اے۔ اینوں ٹھیک طرح تعریف کرن لئی پہلاں اسیں سیٹاں دیاں \emph{نقلاں} لین دا خیال ورتدے آں۔ طریقہ ایہ اے کہ اپنی مرضی دے نشان، $0$ تے $1$، لا کے چیتا رکھیے کہ کیہڑی شے کتھوں آئی۔

\begin{defn}\ollabel{defdissum}
$A$ تے $B$ دا \emph{وکھ وکھ حصیاں والا جوڑ} $A \disjointsum B = (A\times
\{0\}) \cup (B \times \{1\})$ اے۔
\end{defn}

اگے اسیں ترتیبی عدداں دیاں جوڑیاں اُتے اک ترتیب تعریف کردے آں۔

\begin{defn}
کسے وی ترتیبی عدداں $\alpha_1, \alpha_2, \beta_1, \beta_2$ لئی ایویں آکھو:
\begin{align*}
	\tuple{\alpha_1, \alpha_2} \rlexless \tuple{\beta_1, \beta_2}\text{ اُدوں تے صرف اُدوں جدوں }&
	\text{یا تاں $\alpha_2 \in \beta_2$}\\
	& \text{یا دوویں $\alpha_2 = \beta_2$ تے $\alpha_1 \in \beta_1$}
\end{align*}
\end{defn}

ایہ \emph{اُلٹی لغتی ترتیب} اے، کیوں جے پہلاں دوجا رکن ویکھ کے ترتیب لاندے او، تے فیر پہلا۔ چیتے کرو کہ اسیں $\alpha \ordplus \beta$ نوں $\alpha$ دی اک نقل تے اوہدے پچھوں $\beta$ دی اک نقل دی ترتیب دی قسم وجوں تعریف کرنا چاہندے ساں۔ ایہ کرن لئی اسیں آکھدے آں:

\begin{defn}\ollabel{defordplus}
کسے وی ترتیبی عدداں $\alpha$، $\beta$ لئی، اوہناں دا جوڑ $ \alpha \ordplus
\beta = \ordtype{\alpha \disjointsum \beta, \rlexless}$ اے۔
\end{defn}
\noindent
دھیان دیو کہ ایتھے اسیں لکھت نوں ذرا کھلّھ نال ورتیا اے۔ پورا ٹھیک لکھن لئی «$\Setabs{\tuple{x,y}\in \alpha\disjointsum\beta}{x \rlexless y}$» لکھنا چاہیدا اے، «$\rlexless$» دی تھاں۔ پر مختصر رکھن لئی اسیں اگے وی ایویں ای لکھاں گے۔

ہِیٹھلا نتیجہ، \olref[ordinals][ordtype]{thmOrdinalRepresentation} نال رَل کے، پکا کردا اے کہ ساڈی تعریف ٹھیک بنی اے۔

\begin{lem}\ollabel{ordsumlessiswo}
$\tuple{\alpha \disjointsum \beta, \rlexless}$ حُسنِ ترتیب والی ساخت اے، کسے وی ترتیبی عدداں $\alpha$ تے $\beta$ لئی۔
\end{lem}

\begin{proof}
صاف اے کہ $\rlexless$، $\alpha \disjointsum \beta$ اُتے جُڑیا ہویا اے۔ ایہ وکھاؤن لئی کہ اوہ بنیاد والا تعلق اے، کوئی غیر خالی $X \subseteq \alpha
\disjointsum \beta$ مقرر کرو۔ $Y$، $X$ دا اوہ ذیلی سیٹ ہووے جیہدے رکناں دا دوجا محدد جِنّا ہو سکے چھوٹا اے، یعنی $Y = \Setabs{\tuple{\gamma, i} \in X}{(\forall \tuple{\delta, j} \in X)i \leq j}$۔ ہُن $Y$ دا اوہ رکن چُنو جیہدا پہلا محدد سبھ توں چھوٹا اے۔
\end{proof}
\noindent
سو ساڈے کول ترتیبی جوڑ دی اک صاف تے کھلّھی تعریف اے۔ ہِیٹھلی گل کوئی حیرانی والی نہیں (چیتے کرو کہ $1 = \{0\}$، \olref[z][infinity-again]{defnomega} مطابق):
\begin{prop}
$\alpha \ordplus  1 = \ordsucc{\alpha}$، کسے وی ترتیبی عدد $\alpha$ لئی۔
\end{prop}

\begin{proof}
ہم ساخت نقشہ $f$ ویکھو، جیہڑا $\ordsucc{\alpha} = \alpha \cup
\{\alpha\}$ توں $\alpha\disjointsum1 = (\alpha \times \{0\})
\disjointsum (\{0\} \times \{1\})$ ول اے، تے $f(\gamma) =
\tuple{\gamma, 0}$ نال دتا اے $\gamma \in \alpha$ لئی، تے $f(\alpha) = \tuple{0,
1}$۔
\end{proof}
\noindent
نالے وکھانا سوکھا اے کہ جوڑ کجھ بازگشتی شرطاں اُتے پورا اتردا اے۔

\begin{lem}\ollabel{ordadditionrecursion}
کسے وی ترتیبی عدداں $\alpha, \beta$ لئی:
\begin{align*}
	\alpha\ordplus 0 &= \alpha\\
	\alpha \ordplus  (\beta\ordplus 1) &= (\alpha \ordplus  \beta) \ordplus  1\\
	\alpha  \ordplus  \beta &= \supstrict_{\delta < \beta}(\alpha \ordplus  \delta) && \text{جے $\beta $ حدّی ترتیبی عدد اے}
\end{align*}
\end{lem}

\begin{proof}
اسیں واری واری ہر صورت ویکھدے آں۔ پہلاں:
\begin{align*}
	\alpha \ordplus  0 
	& = \ordtype{(\alpha \times \{0\}) \cup (0 \times \{1\}), \rlexless} \\
	&= \ordtype{(\alpha \times \{0\}) \cup \{0\}, \rlexless}\\
	&= \alpha\\
	\alpha \ordplus (\beta \ordplus  1) 
	%&= \ordtype{(\alpha\times \{0\}) \cup (\ordtype{\beta \ordplus  1}\times \{1\}), \rlexless} \\
	&= \ordtype{(\alpha\times \{0\}) \cup (\ordsucc{\beta}\times \{1\}), \rlexless} \\
	&= \ordtype{(\alpha\times \{0\}) \cup (\beta \times \{1\}), \rlexless} \ordplus 1\\
	&= (\alpha \ordplus  \beta) \ordplus  1
\end{align*}
ہُن $\beta \neq \emptyset$ حدّی ہووے۔ جے $\delta < \beta$ تاں $\delta\ordplus 1 < \beta$ وی، سو $\alpha \ordplus  \delta$، $\alpha \ordplus  \beta$ دا حقیقی مُڈھلا ٹکڑا اے۔ سو $\alpha
\ordplus  \beta$، $X = \Setabs{\alpha
\ordplus  \delta}{\delta < \beta}$ دی سخت اُپَرلی حد اے۔ نالے جے $\alpha \leq \gamma <
\alpha \ordplus  \beta$ تاں صاف اے کہ $\gamma = \alpha \ordplus
\delta$ کسے $\delta < \beta$ لئی۔ سو $\alpha \ordplus  \beta =
\supstrict_{\delta< \beta}(\alpha\ordplus \delta)$۔
\end{proof}

پر ایہ اک خاص گل اے۔ ترتیبی جوڑ دی تعریف کرن لئی اسیں \emph{اِس دی تھاں} بس لامتناہی درجیاں والی بازگشت دا قضیہ ورت سکدے ساں، تے بازگشت دیاں مساواتاں این اویویں رکھ سکدے ساں جیویں \olref{ordadditionrecursion} وچ دتیاں نیں (بھاویں «$\ordsucc{\beta}$» ورتدے، «$\beta \ordplus 1$» دی تھاں)۔

سو ترتیبی عدداں اُتے عملاں دی تعریف دے دو وکھرے طریقے نیں۔ اسیں اک حُسنِ ترتیب والا سیٹ صاف صاف بنا کے تے اوہدی ترتیب دی قسم ویکھ کے، \emph{بناوٹ راہیں} عمل دی تعریف کر سکدے آں۔ یا بس بازگشت دیاں مساواتاں رکھ کے، \emph{بازگشت راہیں} تعریف کر سکدے آں۔ پر ٹھیک طرح کریے تاں حاصل اکّو اے۔ کیوں جے \olref[ordinals][ordtype]{thmOrdinalRepresentation} ضمانت دیندا اے کہ ایہ مساواتاں \emph{اکّو اک} ترتیبی عدد طے کردیاں نیں۔

کئی طرحاں ترتیبی حساب، قدرتی عدداں دے جوڑ وانگ ای چلدا اے۔ مثلاً اسیں ایہ ثابت کر سکدے آں:

\begin{lem}\ollabel{ordinaladditionisnice}
جے $\alpha, \beta, \gamma$ ترتیبی عدد نیں، تاں:
\begin{enumerate}
	\item\ollabel{ordaddition1} جے $\beta < \gamma$ تاں $\alpha
	\ordplus  \beta < \alpha \ordplus  \gamma$
	\item\ollabel{ordaddition2} جے $\alpha \ordplus  \beta =
	\alpha\ordplus \gamma$ تاں $\beta = \gamma$
	\item\ollabel{ordaddition3} $\alpha \ordplus  (\beta \ordplus
	\gamma) = (\alpha \ordplus  \beta) \ordplus  \gamma$، یعنی جوڑ شرکتیت رکھدا اے۔
	\item\ollabel{ordaddition4} جے $\alpha \leq \beta$ تاں $\alpha
	\ordplus  \gamma \leq \beta \ordplus \gamma$
\end{enumerate}
\end{lem}

\begin{proof}
اسیں \olref{ordaddition3} ثابت کردے آں تے باقی مشق لئی چھڈدے آں۔ ثبوت $\gamma$ اُتے سادہ لامتناہی درجیاں والے استقرا نال اے، \olref{ordadditionrecursion} ورت کے۔ جدوں $\gamma = 0$:
\[
(\alpha \ordplus  \beta) \ordplus  0 = \alpha \ordplus  \beta  = \alpha \ordplus  (\beta \ordplus  0)
\]
جدوں $\gamma = \delta\ordplus 1$، استقرا لئی منّ لو $(\alpha
\ordplus  \beta) \ordplus  \delta = \alpha \ordplus  (\beta \ordplus
\delta)$۔ ہُن \olref{ordadditionrecursion} تِنّ وار ورت کے:
\begin{align*}
	(\alpha \ordplus  \beta) \ordplus  (\delta \ordplus  1) & = ((\alpha \ordplus  \beta) \ordplus  \delta)\ordplus 1\\
	& = (\alpha \ordplus  (\beta \ordplus  \delta)) \ordplus  1\\
	& = \alpha \ordplus  ((\beta \ordplus  \delta)\ordplus 1)\\
	& = \alpha \ordplus  (\beta \ordplus  (\delta\ordplus 1))
\end{align*}
جدوں $\gamma$ حدّی ترتیبی عدد اے، استقرا لئی منّ لو کہ جے $\delta \in \gamma$ تاں $(\alpha \ordplus  \beta) \ordplus  \delta =
\alpha \ordplus  (\beta \ordplus  \delta)$۔ ہُن:
\begin{align*}
	(\alpha \ordplus  \beta) \ordplus  \gamma & = \supstrict_{\delta < \gamma}((\alpha \ordplus  \beta) \ordplus  \delta) \\
	&= \supstrict_{\delta < \gamma}(\alpha \ordplus  (\beta \ordplus  \delta))\\
	&= \alpha \ordplus  \supstrict_{\delta < \gamma}(\beta \ordplus  \delta)\\
	& = \alpha \ordplus  (\beta \ordplus  \gamma)
\end{align*}
\end{proof}

\begin{prob}
\olref[sth][ord-arithmetic][add]{ordinaladditionisnice} دا باقی حصہ ثابت کرو۔
\end{prob}

ایہناں طرحاں ترتیبی جوڑ تہانوں بڑا جانیا پچھانیا لگنا چاہیدا اے۔ پر اک بہت ضروری گل وچ ترتیبی جوڑ قدرتی عدداں دے جوڑ ورگا \emph{نہیں} اے۔

\begin{prop}\ollabel{ordsumnotcommute}
ترتیبی جوڑ {تبادلی نہیں} اے۔ $1 \ordplus  \omega = \omega <
\omega \ordplus  1$۔
\end{prop}

\begin{proof}
دھیان دیو کہ $1 \ordplus  \omega = \supstrict_{n < \omega} (1 \ordplus n)
= \omega \in \omega \cup \{\omega\} = \ordsucc{\omega} = \omega
\ordplus  1$۔
\end{proof}
\noindent
بھاویں پہلی واری ایہ گل حیران کرے، \emph{ایہ گل حیرانی والی نہیں ہونی چاہیدی}۔ اک پاسے، جدوں تسی $1 \ordplus  \omega$ ویکھدے او، تسی اوہ ترتیب دی قسم سوچدے او جیہڑی سارے قدرتی عدداں توں \emph{پہلاں} اک ہور رکن رکھن نال ملدی اے۔ \olref[sfr][infinite][hilbert]{sec} وچ ہلبرٹ دے ہوٹل والی دلیل وانگ سوچیے، تاں سہج سمجھ مطابق ایس وادھو پہلے رکن نال پوری ترتیب دی قسم نوں کوئی فرق نہیں پینا چاہیدا۔ دوجے پاسے، جدوں تسی $\omega \ordplus  1$ ویکھدے او، تسی اوہ ترتیب دی قسم سوچدے او جیہڑی سارے قدرتی عدداں توں \emph{پچھوں} اک ہور رکن رکھن نال ملدی اے۔ تے ایہ بالکل وکھری گل اے!

\end{document}
